\subsection{Divisor Chooser Continued}

\content{
        \begin{example} Suppose that Betty and Cade are sharing a cake which is
        half strawberry and half vanilla. If there are six slices of cake total,
        and cade values straberry 3 times as much as he values vanilla, how much
        is each slice of straberry worth to him? 
        \end{example}
}{%presentation
\vspace{3in}
}{% nopresentation
\vspace{3in}
}{% comments
}

\content{
        \begin{example} Your turn! If Betty (in the same problem as above)
        values strawberry 2 times as much as vanilla, how much is each slice of
        strawberry worth to betty? 
        \end{example}
}{%presentation
\vspace{3in}
}{% nopresentation
\vspace{3in}
}{% comments
}


\content{
        \begin{example} Two retirees, Pierce and Buzz, purchase a vacation home
        together in Florida, with the agreement that they will split up the year
        into two pieces which will determine when each can use the home. Pierce
        loves Florida in the winter months, and dislikes summer, while Buzz
        likes all parts of the year equally. If Pierce is the divider, describe
        how he might divide up the year into two pieces. Which piece would Buzz
        choose? (e.g. Pierce could do something like the following: May-August
        is one piece, and the other is September-April). There are many correct
        answers to this, explain why yours produces a fair division. 
        \end{example}
}{%presentation
\vspace{4in}
}{% nopresentation
\vspace{4in}
}{% comments
}
